\documentclass[11pt]{article}
\usepackage[margin=25mm]{geometry}
\usepackage{fontspec,unicode-math}
\setmainfont{Latin Modern Roman}
\setmathfont{Latin Modern Math}
\usepackage{amsmath,amssymb,mathtools}
\usepackage{xurl,hyperref}
\hypersetup{colorlinks=true,urlcolor=blue,linkcolor=blue}
\setcounter{secnumdepth}{0}
\setlength{\parindent}{0pt}
\setlength{\parskip}{5pt}
\setlength{\emergencystretch}{3em}
\newcommand{\sref}[2]{\hyperlink{source:#1:#2}{[S#2]}}
\newcommand{\cataloguescope}[1]{\paragraph{Source question.}#1}
\newcommand{\floor}[1]{\left\lfloor#1\right\rfloor}
\begin{document}
\section{Minimal Convex Difference Decompositions of CPWL Functions}\label{statement}

\subsection{Definitions and mathematical statement}

\paragraph{Shared notation.}Standard number systems are denoted by $\mathbb N,\mathbb Z,\mathbb Q,\mathbb R,\mathbb C$; index ranges, computational models and specialized notation follow the cited problem source and the definitions below.


For continuous piecewise-linear $f$, each affine region carries a linear formula. A convex-difference representation is $f=g-h$ with convex piecewise-linear $g,h$; minimizing pieces requires counting the regions according to the source\textquotesingle s convention.

A continuous piecewise-linear function is continuous and agrees with an affine function on each member of a finite polyhedral partition. A difference-of-convex decomposition writes it as $g-h$ with both $g$ and $h$ convex.
\paragraph{Statement recorded in the supplied source.}
Given a continuous piecewise-linear function $f$, decompose it as $f=g-h$ with $g$ and $h$ convex continuous piecewise-linear functions having as few linear pieces as possible. Construct a finite procedure that finds a globally minimal decomposition in arbitrary dimension. \sref{statement}{2}
\paragraph{Scope qualification.}
The objective counting $g$ and $h$ pieces, encoding, and treatment of compatible subdivisions are not specified. The target is global across decompositions, not optimality on one fixed complex.
\paragraph{Source status.}
The supplied catalogue labels this record partially solved. Its source-status metadata was checked against the supplied snapshot on 2026-10-11. This import does not independently establish current open status or mathematical correctness. The catalogue literature assessment is dated 2026-08-22.
\subsection{Short English statement}
Find a finite globally optimal procedure for decomposing a piecewise-linear function as a difference of two convex ones.
\subsection{Sources}
\begin{itemize}
\item[\hypertarget{source:statement:1}{S1}] Ulam AI and the UnsolvedMath Contributors, UnsolvedMath: A Curated Collection of Open Mathematics Problems, catalogue entry. Catalogue curation and metadata: CC BY 4.0. The supplied source question is retained; definitions and the English summary are editorial additions.
\url{https://huggingface.co/datasets/ulamai/UnsolvedMath}.
\item[\hypertarget{source:statement:2}{S2}] Mini-Workshop: Alcoved Polytopes in Physics and Optimization (2025), Oberwolfach Reports, DOI 10.4171/owr/2025/12, pp. 531--558
\url{https://doi.org/10.4171/owr/2025/12}.
\item[\hypertarget{source:statement:3}{S3}] Decomposition polyhedra of CPWL functions
\url{https://openreview.net/forum?id=vVCHWVBsLH}.
\textit{Underlying Oberwolfach Reports material retains its source-specific licence, including CC BY-SA 4.0 where applicable; see THIRD\_PARTY\_NOTICES.md.}
\end{itemize}
\paragraph{Preparation and review.}
Original question attributed to the cited source and UnsolvedMath contributors. Codex (GPT-6; exact serving revision not exposed) assisted with definitions, English summaries, source inspection and formatting. No new human mathematical authorship or independent proof verification is asserted. Catalogue identity checked against the complete supplied registry and existing repository statements; known flagged repetitions and obvious identical conjectures excluded. Full mathematical-equivalence/literature review has not been performed. Statement preparation and final review followed contribution prompts C and D; automated checks verify structure and identity, not open status.
\label{end:statement}
\end{document}
